% =========================================================
% TikZ diagram: Derived Set
% Place immediately after the Derived Set definition box.
% Requires in preamble: \usepackage{tikz}
%                       \usetikzlibrary{calc}
% =========================================================

\begin{figure}[h]
\centering
\begin{tikzpicture}[
  dot/.style     = {circle, fill, inner sep=1.8pt},
  opendot/.style = {circle, draw, fill=white, inner sep=1.8pt, line width=0.8pt},
  lbl/.style     = {font=\small},
  sublbl/.style  = {font=\scriptsize, text=gray!55!black}
]

% --- E (blue blob) ---
\fill[blue!12, draw=blue!50, line width=0.9pt]
  (-0.1, -0.5)
  .. controls ( 0.6, -0.8) and ( 1.4,  0.0) .. ( 1.3,  1.0)
  .. controls ( 1.2,  2.0) and ( 0.8,  3.0) .. (-0.2,  3.3)
  .. controls (-1.0,  3.6) and (-2.6,  3.4) .. (-3.2,  2.4)
  .. controls (-3.8,  1.4) and (-3.6,  0.2) .. (-2.8, -0.4)
  .. controls (-2.0, -1.0) and (-0.8, -0.8) .. (-0.1, -0.5)
  -- cycle;

\node[lbl, text=blue!60!black] at (-1.6, 0.5) {$E$};

% --- E' (orange dashed contour) ---
% Right lune (x > 0.9) lies in E' but not E  → E' \ E
% Top-left pocket (x < -2.5, y > 2.2) lies in E but not E' → E \ E'
\draw[orange!75!black, dashed, line width=1.2pt]
  ( 1.3,  0.0)
  .. controls ( 2.0,  0.6) and ( 2.0,  2.6) .. ( 1.1,  3.5)
  .. controls ( 0.2,  4.2) and (-1.4,  3.8) .. (-2.2,  3.0)
  .. controls (-2.8,  2.4) and (-2.6,  1.0) .. (-2.2,  0.2)
  .. controls (-1.8, -0.6) and ( 0.2, -0.8) .. ( 1.3,  0.0)
  -- cycle;

\node[lbl, text=orange!75!black] at (2.55, 1.8) {$E'$};

% --- Ambient space ---
\node[lbl, font=\small\itshape] at (-4.4, 3.8) {$X$};

% --- Region labels ---
\node[sublbl] at (-0.6, -0.1) {$E \cap E'$};
\node[sublbl, align=center] at ( 1.8,  0.6) {$E'{\setminus}E$};

% --- a in E ∩ E' : non-isolated interior point ---
\node[dot, fill=blue!70!black] (A) at (-1.0, 1.6) {};
\node[lbl, anchor=south west] at ($(A)+(0.10, 0.10)$) {$a$};
\node[sublbl, anchor=north west] at ($(A)+(0.10,-0.10)$)
  {$a\in E,\;a\in E'$};

% --- b in E' \ E : limit point of E not in E ---
\node[opendot, draw=orange!75!black] (B) at (1.05, 1.8) {};
\node[lbl, anchor=south west] at ($(B)+(0.10, 0.10)$) {$b$};
\node[sublbl, anchor=north west] at ($(B)+(0.10,-0.10)$)
  {$b\notin E,\;b\in E'$};

% --- c in E \ E' : isolated point of E ---
\coordinate (C) at (-3.0, 2.6);
\draw[gray!40, dotted, line width=0.8pt] (C) circle (0.48cm);
\node[dot, fill=blue!70!black] at (C) {};
\node[lbl, anchor=east] at ($(C)+(-0.55, 0.13)$) {$c$};
\node[sublbl, anchor=north east, align=right] at ($(C)+(-0.55,-0.13)$)
  {$c\in E,\;c\notin E'$\\[-1pt](isolated)};

% --- Legend ---
\begin{scope}[shift={(-2.6, -1.6)}]
  \fill[blue!12, draw=blue!50, line width=0.7pt] (0,0) rectangle (0.48,0.30);
  \node[lbl, anchor=west] at (0.62, 0.15) {$E$};

  \draw[orange!75!black, dashed, line width=1.2pt]
    (1.6, 0.15) -- (2.2, 0.15);
  \node[lbl, anchor=west] at (2.35, 0.15) {boundary of $E'$};

  \draw[gray!40, dotted, line width=0.8pt]
    (5.4, 0.15) -- (6.0, 0.15);
  \node[lbl, anchor=west] at (6.15, 0.15) {isolation ball around $c$};
\end{scope}

\end{tikzpicture}

\smallskip
{\small
  $a \in E \cap E'$: a non-isolated point of $E$, and a limit point of $E$.
  \quad
  $b \in E' \setminus E$: a limit point of $E$ not belonging to $E$.
  \quad
  $c \in E \setminus E'$: an isolated point of $E$, so $c \notin E'$.
}
\caption{The derived set $E' = \{x \in X : x \text{ is a limit point of } E\}$
relative to $E$.
The three annotated points illustrate the meaningful membership
combinations for points of~$X$.}
\label{fig:derived-set}
\end{figure}